\section{The energy shocks and the adoption channel}
\label{sec:experiments}

\subsection{The two shocks and the common-steady-state counterfactual}
\label{sec:counterfactual}

The \emph{price} shock is a 100 log-point rise in the world energy price with a
16-quarter half-life, as in \citet{ARS2024}: energy supply is perfectly elastic
($\varepsilon^{E}=\infty$) and the world price is exogenous. The \emph{supply}
shock is a 10\% cut in available energy for 6 quarters, after which the world
price clears a fixed quantity endogenously, as in \citet{Bayer2026}. The two are
easiest to read side by side (Figure~\ref{fig:shocks}): under the price shock we
move the world price directly and the quantity clears; under the supply shock we
move the quantity, a 10\% cut, and the world price clears it. Because energy
demand is highly inelastic ($\eta_E=0.10$), a 10\% quantity cut requires a price
spike of about $170\%$ to clear, so the supply shock is in fact the larger
\emph{price} disturbance of the two on impact. Impulse
responses are reported as $100\times$ the level deviation from steady state.
Variables whose steady state is one (prices, $w$, $n$, and the shares $D^{G}$,
$D^{sw}$) read directly as percent or percentage-point deviations; $C$ (steady
state $\approx0.997$) is within rounding of a percent; $y$ and $\mathrm{gdp}$
(steady state $\approx0.985$) are level deviations $\times100$, within about
$1.5\%$ of the exact percentage; energy quantities are a share of the energy
steady state ($\approx0.040$).

Several results below isolate the adoption channel as the difference between the
full model and a ``no-adoption'' counterfactual,
$\Delta_H\equiv\sum_{t<H}(y^{\mathrm{adopt}}_t-y^{\mathrm{no\text{-}adopt}}_t)$.
The counterfactual is constructed to share the full model's steady state
\emph{exactly}: households still hold the discrete durable-choice margin, and the
steady-state green share, switching flow, and switching cost
($\bar D^{G}=5\%$, $D^{sw}_{ss}$, $\psi_g$) are bit-identical to the adoption
economy. What differs is the \emph{transition} only: the discrete-choice
probabilities are held at their steady-state values throughout the impulse
response, so the durable composition does not respond to the crisis while every
other block responds normally. We call this the \textbf{common-steady-state}, or
\emph{frozen-choice}, counterfactual. The alternative, shutting the choice by
making green durables infeasible, so the counterfactual's own steady state has
$D^{G}_{ss}=0$,would compare two economies that differ both in the crisis
response \emph{and} in steady-state composition; the frozen construction isolates
the crisis response alone.

\begin{figure}[H]\centering
\IfFileExists{output/fig_shocks_import.pdf}{\includegraphics[width=\textwidth]{output/fig_shocks_import.pdf}}{\textit{[figure output/fig\_shocks\_import.pdf: run run\_shocks.py]}}
\caption{The two energy shocks. \emph{Left}: the price shock moves the world energy
price (exogenous) and the quantity clears. \emph{Right}: the supply shock moves
availability ($-10\%$ for 6 quarters, exogenous) and the world price clears the
reduced quantity, spiking about $170\%$ because demand is inelastic. In the right
panel the bold blue line is the announced availability cut and the dashed blue line
is realized availability: the two differ because world stock arbitrage
\eqref{eq:energy-supply} releases inventories against the cut, buffering part of it,
and rebuilds them afterwards. Percent deviations from own steady state.}
\label{fig:shocks}
\end{figure}

Figures~\ref{fig:baseline} and~\ref{fig:baseline_supply} report the baseline
responses to the two shocks, each with the adoption margin open (solid) against the
frozen counterfactual (dashed).

\begin{figure}[H]\centering
\includegraphics[width=0.92\textwidth]{output/fig1_price_shock.png}
\caption{Brown energy \emph{price} shock (elastic supply), adoption open (solid)
vs.\ frozen at the common steady state (dashed), no policy.}
\label{fig:baseline}
\end{figure}

\begin{figure}[H]\centering
\includegraphics[width=0.92\textwidth]{output/fig1_supply_shock.png}
\caption{Brown energy \emph{supply} shock ($-10\%$ for 6 quarters, fixed
quantity), adoption open (solid) vs.\ frozen (dashed), no policy.}
\label{fig:baseline_supply}
\end{figure}

\subsection{What the adoption margin adds}
\label{sec:signmap}\label{sec:crisis}

Under the import booking and the price shock, the adoption channel's contribution
to cumulative output is $\Delta_{24}=-8.8$ (Table~\ref{tab:summary};
Figure~\ref{fig:decomposition} plots the margin's contribution across variables): the margin
is \emph{contractionary}. It follows a race between two resource flows in the
balance of payments (Section~\ref{sec:bop}). Adoption imports a \textbf{switching
cost} $\psi_g D^{sw}_t$, a resource outflow, \emph{front-loaded} (paid once by
each switcher) and scaling with the adoption \emph{volume} $D^{sw}$,and, against
it, a \textbf{flow energy saving} $(P^{E}_{B}-P^{E}_{G})\,C^{G}_{E,t}$ booked as
reduced imports, an inflow accruing \emph{gradually} over the life of the green
durable. Under the frozen counterfactual and the import booking the front-loaded
cost dominates the windowed saving at every persistence. The magnitude is the
pocket cost of the crisis-time switching wave at the calibrated $\psi_g=0.54$
($\approx0.54\times C_{ss}$), not a general-equilibrium artefact: the direct cost
$\psi_g\times$ extra switchers cumulates to $+8.96$, matching the measured gap
almost exactly. Under the fixed-quantity (supply) closure a third, expansionary
term appears, adoption lowers fossil demand against a fixed supply, so the
clearing price falls and relaxes scarcity for \emph{all} households, which shifts
$\Delta_{24}$ up to $-0.35$. The marginal switcher's steady-state indifference,
$\psi_g\approx(\bar P^{E}_{B}-\bar P^{E}_{G})\,\bar C^{G}_{E}/(r+\delta_g)$, ties
the threshold to the cost ratio $\varrho$.

\paragraph{Booking caveat.} The \emph{sign} of $\Delta_{24}$, not the policy
ordering below, depends on how the green resource flows are booked. Under the
alternative domestic-industry booking (green energy and installation as
zero-profit domestic value added rather than imports) it flips to $+28.6$ under
the price shock and $+4.7$ under the supply shock. We report the import booking
throughout and collect the comparison in Appendix~\ref{app:booking}. This is a
booking convention, not a result about the transition, and it is why we phrase the
paper's headline in terms of the \emph{margin's response} to policy rather than
its output sign.

\begin{figure}[H]\centering
\includegraphics[width=\textwidth]{output/fig5_decomposition.png}
\caption{Contribution of the green adoption margin (full minus frozen), by
shock. Output, consumption, energy use, net exports.}
\label{fig:decomposition}
\end{figure}

The race between the two resource flows is exact on the import side.
Figure~\ref{fig:adoption_flows} splits the adoption channel's contribution to
imports, full minus frozen, under the baseline import booking, into the
front-loaded switching-cost outflow $\psi_g D^{sw}_t$ and everything else, the
latter dominated by the gradual energy saving from the brown-to-green mix shift.
The split is exact by construction. The switching cost dominates the saving at
every horizon under both shocks: under the price shock it cumulates to $+8.96$
against a $-5.63$ energy-mix inflow, a net import response of $+3.33$ that accounts
for the channel's contractionary output contribution ($\Delta_{24}=-8.8$). The
outflow is largest exactly on impact, when household income is falling hardest,
which is why booking it as an import makes the channel contractionary rather than
its timing incidental.

\begin{figure}[H]\centering
\includegraphics[width=\textwidth]{output/fig5b_adoption_flows.png}
\caption{Adoption channel on the import side (full minus frozen): the front-loaded
switching-cost outflow $\psi_g D^{sw}$ against the gradual energy-saving inflow, and
their sum. Contributions in level deviations $\times100$; import booking.}
\label{fig:adoption_flows}
\end{figure}

The adoption channel's output effect is a domestic-demand phenomenon. Decomposing
output one level below the domestic/export split (Figure~\ref{fig:output_channels}),
the frozen-choice counterfactual separates domestic demand into what it would be with
the margin shut, the AMRS real-income and substitution channels, and the adoption
channel's increment. The adoption channel's cumulative contribution ($-9.0$) falls
almost entirely on domestic home-good demand, not exports ($+0.2$): the
switching-cost import outflow draws resources out of domestic absorption. Its
magnitude nearly matches the export offset ($+8.7$) with the opposite sign, so
aggregate output tracks the AMRS channel while the adoption drag and the
competitiveness gain roughly cancel.

\begin{figure}[H]\centering
\IfFileExists{output/fig_output_channels_import.pdf}{\includegraphics[width=0.8\textwidth]{output/fig_output_channels_import.pdf}}{\textit{[figure output/fig\_output\_channels\_import.pdf: run run\_output\_channels.py]}}
\caption{Output response, three exact channels: exports $c_H^{*}$, domestic demand
with the adoption margin frozen (AMRS), and the adoption channel's increment to
domestic demand. The three sum to $y$. Level deviations $\times100$, price shock, no
policy.}
\label{fig:output_channels}
\end{figure}

The same response, read from the household budget rather than from output, is almost
entirely a general-equilibrium phenomenon. Decomposing consumption through the
household consumption Jacobians, $dC=\sum_x J^{C,x}\,dx$ over the household's inputs,
exact to first order, separates the direct effect of the shock, the energy bill in
the budget, from the indirect effects that run through factor income and the real
rate (Figure~\ref{fig:direct_indirect}, computed in the CPI numeraire where these
inputs are unambiguous; the total is numeraire-invariant). The direct energy-price
channel is negligible ($-0.2$ cumulative): energy is a small budget share, so a
higher energy price does little to consumption on its own. The collapse is the
indirect real-income channel ($-101.8$ of a $-105.6$ total), with the real rate
adding $-3.6$. The energy shock reaches consumption almost entirely by lowering
factor income in general equilibrium, not by raising the household's energy bill, 
the mechanism \citet{ARS2024} and the broader HANK literature identify, here in an
open economy with an adoption margin.

\begin{figure}[H]\centering
\IfFileExists{output/fig_direct_indirect_import.pdf}{\includegraphics[width=0.78\textwidth]{output/fig_direct_indirect_import.pdf}}{\textit{[figure output/fig\_direct\_indirect\_import.pdf: run run\_direct\_indirect.py]}}
\caption{Consumption response to the brown-price shock, decomposed into the channels
through which the shock reaches the household budget: the direct energy-price effect
and the indirect effects operating through real income and the real rate. Household
consumption Jacobians, CPI numeraire; the channels sum to the total. Level deviations
$\times100$.}
\label{fig:direct_indirect}
\end{figure}

\paragraph{Who the shock recruits.} The crisis draws in households who would not have
switched at the pre-crisis steady state. Figure~\ref{fig:induced_adopters}
characterises them with a comparative static: the switching probability of the median
type at a steady state with a permanently $10\%$ higher brown price, minus the
baseline field. The induced switching is positive throughout and rises with
productivity; at high productivity the recruitment reaches down to more moderate
wealth, pulling in households who sat just below the switching threshold. Weighted by
the steady-state distribution, the induced flow concentrates at moderate productivity
and low-to-moderate wealth, where the brown mass sits. The construction is a
comparative-static proxy: it uses a permanent price change rather than the transitory
crisis path, so it locates and signs the induced margin rather than timing it; the
exact transition object reads the disaggregated choice probabilities along the impulse
response.

\begin{figure}[H]\centering
\IfFileExists{output/fig_induced_adopters_import.pdf}{\includegraphics[width=0.8\textwidth]{output/fig_induced_adopters_import.pdf}}{\textit{[figure output/fig\_induced\_adopters\_import.pdf: run run\_induced\_adopters.py]}}
\caption{Induced adopters: the extra quarterly switching probability a permanently
$10\%$ higher brown price induces, median discount-factor type, over the
wealth--productivity grid (comparative static). Contours mark steady-state brown mass.}
\label{fig:induced_adopters}
\end{figure}

\input{output/tab_summary_core_import.tex}
